\chapter[1936 --- Dale et al.]{1936: Henry Hallett Dale, Otto Loewi}
\label{ch:1936_Dale}

\section*{Main Contribution}

\textit{Discovered that nerve impulses are transmitted chemically (not electrically) across synapses, using acetylcholine as the neurotransmitter---founding the field of neuropharmacology.}

\section{Official Citation}

for their discoveries relating to chemical transmission of nerve impulses

\section{Background and Historical Context}

\subsection{Henry Hallett Dale}

Henry Hallett Dale was born on June 9, 1875, in London, England. He studied at Trinity College, Cambridge, where he excelled in the natural sciences, and later at St. Bartholomew's Hospital Medical College, where he qualified in medicine in 1900. Dale's early career was marked by a broad range of interests: he studied the pharmacology of amines and alkaloids at the University of London, worked at the Pasteur Institute in Paris, and served as director of the Wellcome Physiological Research Laboratories from 1914 to 1928. In 1928, Dale was appointed director of the National Institute for Medical Research in London, a position he held until his retirement in 1942.

Dale's research focused on the pharmacology of the nervous system---the study of how chemical substances (both naturally occurring and synthetic) affect the function of nerve cells and their target tissues. His early work on histamine, ergot alkaloids, and the effects of organophosphorus compounds on the nervous system established him as one of the leading pharmacologists of his generation, and it laid the groundwork for his Nobel Prize-winning work on the chemical transmission of nerve impulses. Dale's research was characterized by meticulous experimental technique, rigorous logical analysis, and a deep understanding of the pharmacological properties of the substances he studied.

\subsection{Otto Loewi}

Otto Loewi was born on June 3, 1873, in Frankfurt, Germany. He studied medicine at the University of Strassburg, where he received his medical degree in 1896, and later at the universities of Munich and Wiesbaden. Loewi's research career began in the laboratory of the pharmacologist Hans Horst Meyer in Marburg, where he studied the effects of various drugs on the heart. In 1905, Loewi was appointed professor of pharmacology at the University of Graz in Austria, a position he held until 1938, when he was forced to flee the Nazi annexation of Austria (the Anschluss). Loewi was arrested by the Gestapo in 1938, and he was eventually released and allowed to emigrate to the United States in 1940, where he held a research position at the New York University School of Medicine until his death in 1961.

Loewi's research was characterized by elegant simplicity and a deep appreciation of the experimental method. His Nobel Prize-winning experiment---in which he demonstrated that the vagus nerve released a chemical substance that slowed the heartbeat---is one of the most famous experiments in the history of physiology, and it established the principle of chemical neurotransmission that underpins our entire understanding of how the nervous system communicates with its target tissues.

\section{The Discovery and Contribution}

\subsection{The Historical Background: Chemical vs. Electrical Transmission}

The question of how nerve impulses are transmitted across the gap between nerve cells (the synapse) and between nerve cells and their target tissues (the neuromuscular junction and the neuroeffector junction) was one of the most contentious issues in physiology in the late nineteenth and early twentieth centuries. The two competing hypotheses were the ``electrical'' theory, which held that nerve impulses were transmitted by the local flow of electrical current, and the ``chemical'' theory, which held that nerve impulses were transmitted by the release of chemical substances from nerve terminals.

The electrical theory had a distinguished pedigree, dating back to the work of Emil du Bois-Reymond in the 1840s. It was supported by the observation that electrical stimulation of nerve fibers produced rapid, virtually instantaneous responses in their target tissues, and it was consistent with the known properties of electrical conduction in nerve fibers. The chemical theory, while less widely accepted, was supported by the observation that certain drugs (such as adrenaline and acetylcholine) could mimic the effects of nerve stimulation, and that the effects of nerve stimulation could be blocked by specific chemical antagonists. Thomas R. Elliott, a young British physician, had proposed in 1905 that the sympathetic nervous system might transmit its impulses by releasing adrenaline-like substances, but his hypothesis had not been tested experimentally.

The debate between the electrical and chemical theories was not merely academic: it had significant implications for the understanding of nervous system function and for the development of pharmacological therapies. If neurotransmission was electrical, then pharmacological agents would have little to no effect on synaptic transmission, and the treatment of neurological and psychiatric disorders would have to rely on electrical stimulation or surgical intervention. If neurotransmission was chemical, then pharmacological agents could be used to modulate synaptic transmission, opening up an enormous range of therapeutic possibilities.

\subsection{Loewi's Experiment}

Loewi's experiment, which provided the first direct evidence for chemical neurotransmission, was conducted on the night of March 31, 1921, and it has become one of the most celebrated episodes in the history of science. The story, as Loewi later recounted it, began with a dream: Loewi dreamed of an experiment that would demonstrate the chemical nature of vagus nerve transmission, and he woke up in the middle of the night, wrote down the essential details of the experiment, and went back to sleep. The next morning, he could not read his handwriting, but the experiment was so vivid in his memory that he went to the laboratory and performed it immediately.

The experiment was deceptively simple. Loewi dissected two frog hearts: one with the vagus nerve still attached (the ``donor'' heart) and one without (the ``recipient'' heart). He stimulated the vagus nerve of the donor heart electrically and observed that the heart rate slowed down (as was well known---the vagus nerve was the ``inhibitory'' nerve of the heart). He then took the solution that had been perfusing the donor heart and transferred it to the recipient heart, and he observed that the recipient heart also slowed down, even though its own vagus nerve had not been stimulated.

This result was significant: it demonstrated that the vagus nerve, when stimulated, released a chemical substance into the perfusing solution, and that this substance was capable of slowing the heartbeat when applied to another heart. The experiment ruled out the electrical theory of neurotransmission, because the electrical current from the vagus nerve could not have been transferred in the perfusing solution. It proved that the vagus nerve communicated with the heart by releasing a chemical substance---what Loewi called ``Vagusstoff'' (vagus substance)---and that this substance was the mediator of the inhibitory effect of the vagus nerve on the heart.

Loewi repeated the experiment many times, and the results were consistently reproducible. He also showed that the effects of the Vagusstoff were potentiated by eserine (an inhibitor of the enzyme that degrades the substance), and that the substance was destroyed by boiling---properties that were consistent with a small organic molecule. Loewi published his findings in 1921, and the paper was immediately recognized as one of the major in the history of physiology.

Loewi's Vagusstoff was later identified as acetylcholine (ACh), the first neurotransmitter to be discovered. The identification was achieved by Dale and his colleagues, who showed that the effects of Vagusstoff were identical to those of acetylcholine in every respect: it was hydrolyzed by the enzyme acetylcholinesterase, it was potentiated by eserine (an acetylcholinesterase inhibitor), and it produced the same pattern of responses as acetylcholine on a variety of smooth muscle preparations. The identification of Vagusstoff as acetylcholine was a triumph of pharmacological analysis, and it established the principle that nerve cells communicate by releasing chemical neurotransmitters.

\subsection{Dale's Contributions}

Dale's contribution to the understanding of chemical neurotransmission extended far beyond the identification of acetylcholine. He conducted a systematic pharmacological analysis of the effects of acetylcholine and other autonomic neurotransmitters, characterizing their actions on different tissues and defining the concept of ``cholinergic'' and ``adrenergic'' transmission.

Dale showed that the parasympathetic nervous system communicated with its target tissues by releasing acetylcholine (the ``cholinergic'' neurotransmitter), while the sympathetic nervous system communicated by releasing adrenaline or noradrenaline (the ``adrenergic'' neurotransmitters). He also characterized the muscarinic and nicotinic effects of acetylcholine---the distinction between the actions of acetylcholine on smooth muscle (mediated by muscarinic receptors, named because they were activated by muscarine, a toxin from the mushroom \emph{Amanita muscaria}) and on skeletal muscle (mediated by nicotinic receptors, named because they were activated by nicotine)---a distinction that remains fundamental to pharmacology and to the understanding of autonomic nervous system function.

Dale's work on the pharmacology of neurotransmitters established the framework for the entire field of autonomic pharmacology, which is concerned with the drugs that affect the sympathetic and parasympathetic nervous systems. The development of drugs for the treatment of hypertension, asthma, glaucoma, depression, and many other conditions is based on the pharmacological principles that Dale elucidated.

\subsection{The Joint Nobel Prize}

Dale and Loewi were jointly awarded the Nobel Prize in Physiology or Medicine in 1936 ``for their discoveries relating to chemical transmission of nerve impulses.'' The award recognized Loewi's experimental demonstration that the vagus nerve released a chemical substance and Dale's pharmacological characterization of that substance and its relatives. Together, their work established the principle of chemical neurotransmission---the idea that nerve cells communicate by releasing chemical substances---which is now recognized as the fundamental mechanism of synaptic communication in the nervous system.

\section{Impact on Modern Medicine}

The principle of chemical neurotransmission that Dale and Loewi established is the foundation of modern neuropharmacology, which is concerned with the development of drugs that modulate the function of the nervous system by targeting neurotransmitter systems. The vast majority of drugs used in the treatment of neurological, psychiatric, and many other medical conditions act by modifying the synthesis, release, reuptake, or degradation of neurotransmitters, or by binding to neurotransmitter receptors.

In neurology, the understanding of acetylcholine neurotransmission has led to the development of drugs for the treatment of Alzheimer's disease. Cholinesterase inhibitors (such as donepezil, rivastigmine, and galantamine), which increase the availability of acetylcholine in the synaptic cleft by inhibiting the enzyme that breaks it down, are used to improve cognitive function in patients with mild to moderate Alzheimer's disease. These drugs are a direct application of Dale's insights into the cholinergic system, and they represent one of the few pharmacological interventions available for this devastating neurodegenerative disease.

In psychiatry, the understanding of monoamine neurotransmission (serotonin, dopamine, and noradrenaline) has led to the development of a wide range of antidepressants and antipsychotic drugs. Selective serotonin reuptake inhibitors (SSRIs, such as fluoxetine and sertraline), which increase the availability of serotonin in the synaptic cleft by blocking its reuptake into the presynaptic neuron, are among the most widely prescribed drugs in the world for the treatment of depression and anxiety disorders. Antipsychotic drugs, which block dopamine receptors in the brain, are used to treat schizophrenia and other psychotic disorders. These drugs are all descendants of the pharmacological principles that Dale and Loewi established.

In anesthesiology, the understanding of neuromuscular transmission has led to the development of neuromuscular blocking agents (such as succinylcholine and vecuronium) that are used to produce muscle relaxation during surgery. These drugs act by blocking the nicotinic acetylcholine receptors at the neuromuscular junction, a mechanism that was first characterized by Dale.

In cardiology, the understanding of cholinergic and adrenergic neurotransmission has led to the development of drugs for the treatment of arrhythmias, heart failure, and hypertension. Beta-adrenergic blockers (such as propranolol and metoprolol), which block the effects of noradrenaline on the heart, are used to reduce heart rate and blood pressure, while cholinesterase inhibitors are used to treat bradycardia (slow heart rate). These drugs are direct applications of the autonomic pharmacology that Dale established.

The concept of chemical neurotransmission has also been extended beyond the nervous system to the understanding of endocrine signaling, paracrine signaling, and autocrine signaling. The recognition that cells communicate by releasing chemical messengers---whether neurotransmitters, hormones, cytokines, or growth factors---is a general principle of cell biology that has its origins in the work of Dale and Loewi. The identification of over 100 neurotransmitters and neuromodulators, the characterization of their receptors and signaling pathways, and the development of drugs that target these systems are all part of the legacy that Dale and Loewi established.

\section{Legacy}

The legacies of Dale and Loewi are inseparable from the history of modern neuroscience and pharmacology. Loewi, whose elegant experiment demonstrated the chemical nature of neurotransmission, exemplifies the power of simple, well-designed experiments to resolve fundamental questions about the nature of biological systems. His famous ``Vagusstoff'' experiment, performed on a single night in 1921, remains one of the most instructive examples of experimental design in the history of science. The story of the experiment---the dream, the hastily written notes, the early morning in the laboratory---has become a part of the folklore of science, a reminder that inspiration and insight can come at unexpected moments.

Dale, whose systematic pharmacological analysis of neurotransmitters established the framework for modern neuropharmacology, was a scientist of notable breadth and rigor. His classification of cholinergic and adrenergic transmission, his characterization of muscarinic and nicotinic receptors, and his systematic analysis of the effects of neurotransmitter agonists and antagonists set standards for pharmacological research that remain exemplary. Dale's Nobel lecture, in which he summarized the evidence for chemical neurotransmission with characteristic clarity and precision, is still studied as a model of scientific exposition. Dale was knighted in 1932 and raised to the peerage as Baron Dale of Lister in 1943, and he served as president of the Royal Society from 1940 to 1945.

Together, Dale and Loewi established the paradigm of chemical neurotransmission that underpins our entire understanding of how the nervous system functions. Their work demonstrated that the nervous system is not a simple electrical circuit but a complex chemical signaling network, and that the manipulation of this network by pharmacological agents is the basis of modern neurological and psychiatric therapeutics. The legacy of Dale and Loewi is written in every prescription for an antidepressant, every dose of an anesthetic, and every drug that modulates the function of the nervous system. Their discovery that nerve cells speak to each other in the language of chemistry opened up an entirely new era in medicine, and its influence continues to grow as new neurotransmitter systems, receptor subtypes, and signaling pathways are discovered. Loewi died in New York on December 25, 1961, and Dale died in London on July 23, 1968, but their influence on medicine and pharmacology endures as one of the great achievements of twentieth-century science.


\section{Modern Developments}

Dale and Loewi's neurotransmitter discovery has led to modern neuropharmacology. Understanding of cholinergic transmission has enabled treatments for Alzheimer's disease (donepezil, rivastigmine) and myasthenia gravis (pyridostigmine). Adrenergic transmission understanding has led to beta-blockers, alpha-blockers, and bronchodilators. The development of SSRIs for depression, based on serotonin neurotransmission, has transformed psychiatric medicine. Optogenetic control of neurotransmitter release enables precise circuit manipulation.


\section{Bibliography}

\begin{thebibliography}{9}
\bibitem{nobel-1936}
Nobel Prize Outreach, ``The Nobel Prize in Physiology or Medicine 1936,''\emph{NobelPrize.org}.\\
\url{https://www.nobelprize.org/prizes/medicine/1936/summary/}

\bibitem{lecture-1936}
Henry Hallett Dale, ``Nobel Lecture,''\emph{NobelPrize.org}. Winner-authored primary source; downloadable from the official Nobel Prize archive.\\
\url{https://www.nobelprize.org/prizes/medicine/1936/dale/lecture/}
\end{thebibliography}
